\section{AIME I 2004}
        \begin{enumerate}
        	\item (\textit{AIME I 2004}) The digits of a positive integer $n$ are four consecutive integers in decreasing order when read from left to right. What is the sum of the possible remainders when $n$ is divided by 37?


 
 \textbf{Answer: }217



\textbf{Solution 1}

A brute-force solution to this question is fairly quick, but we'll try something slightly more clever:  our numbers have the form ${\underline{(n+3)}}\,{\underline{(n+2)}}\,{\underline{( n+1)}}\,{\underline {(n)}}$$= 1000(n + 3) + 100(n + 2) + 10(n + 1) + n = 3210 + 1111n$, for $n \in \lbrace0, 1, 2, 3, 4, 5, 6\rbrace$.

 Now, note that $3\cdot 37 = 111$ so $30 \cdot 37 = 1110$, and $90 \cdot 37 = 3330$ so $87 \cdot 37 = 3219$. So the \href{https://artofproblemsolving.com/wiki/index.php?title=Remainder}{remainders} are all congruent to $n - 9 \pmod{37}$. However, these numbers are negative for our choices of $n$, so in fact the remainders must equal $n + 28$.

 Adding these numbers up, we get $(0 + 1 + 2 + 3 + 4 + 5 + 6) + 7\cdot28 = \boxed{217}$




\textbf{Solution 2}

There are only $7$ possible values of $n$: $3210, 4321, 5432, \cdots, 9876$. 

 $3210$ gives a remainder of $28$ when divided by $37$. To calculate the remainders of the other integers, notice that each number is $1111$ more than the previous number. Since $1111$ gives a remainder of $1$ when divided by $37$, the remainders of the other integers will be $29, 30, 31, \cdots, 34$. 

 Therefore, our answer is $28 + 29 + 30 + 31 + 32 + 33 + 34 = \frac{28 + 34}{2} \cdot 7 = \boxed{217}$.~Viliciri

 


	\item (\textit{AIME I 2004}) Set $A$ consists of $m$ consecutive integers whose sum is $2m,$ and set $B$ consists of $2m$ consecutive integers whose sum is $m.$ The absolute value of the difference between the greatest element of $A$ and the greatest element of $B$ is 99. Find $m.$

 \textbf{Answer: }201



\textbf{Solution 1}

Note that since set $A$ has $m$ consecutive integers that sum to $2m$, the middle integer (i.e., the median) must be $2$.  Therefore, the largest element in $A$ is $2 + \frac{m-1}{2}$. 

 Further, we see that the median of set $B$ is $0.5$, which means that the "middle two" integers of set $B$ are $0$ and $1$.  Therefore, the largest element in $B$ is $1 + \frac{2m-2}{2} = m$.  $2 + \frac{m-1}{2} > m$ if $m < 3$, which is clearly not possible, thus $2 + \frac{m-1}{2} < m$.

 Solving, we get 
 \begin{align*} m - 2 - \frac{m-1}{2} &= 99\\ m-\frac{m}{2}+\frac{1}{2}&=101\\ \frac{m}{2}&=100\frac{1}{2}.\\ m &= \boxed{201}\end{align*}





\textbf{Solution 2}

Let us give the \href{https://artofproblemsolving.com/wiki/index.php?title=Element}{elements} of our sets names:$A = \{x, x + 1, x + 2, \ldots, x + m - 1\}$ and $B = \{y, y + 1, \ldots, y + 2m - 1\}$.  So we are given that
 \[2m = x + (x + 1) + \ldots + (x + m - 1) = mx + (1 + 2 + \ldots + (m - 1)) = mx + \frac{m(m -1)}2,\] so $2 = x + \frac{m - 1}2$ and $x + (m - 1) = \frac{m + 3}2$ (this is because $x = 2 - \frac{m-1}{2}$ so plugging this into $x+(m-1)$ yields $\frac{m+3}{2}$). Also,
 \[m = y + (y + 1) + \ldots + (y + 2m - 1) = 2my + \frac{2m(2m - 1)}2,\]so $1 = 2y + (2m - 1)$ so $2m = 2(y + 2m - 1)$ and $m = y + 2m - 1$.

 Then by the given, $99 = |(x + m - 1) - (y + 2m - 1)| = \left|\frac{m + 3}2 - m\right| = \left|\frac{m - 3}2\right|$.  $m$ is a \href{https://artofproblemsolving.com/wiki/index.php?title=Positive\_integer}{positive integer} so we must have $99 = \frac{m - 3}2$ and so $m = \boxed{201}$.

 



\textbf{Solution 3}

The thing about this problem is, you have some "choices" that you can make freely when you get to a certain point, and these choices won't affect the accuracy of the solution, but will make things a lot easier for us.

 First, we note that for set $A$

 \[\frac{m(f + l)}{2} = 2m\]

 Where $f$ and $l$ represent the first and last terms of $A$.  This comes from the sum of an arithmetic sequence.

 Solving for $f+l$, we find the sum of the two terms is $4$.

 Doing the same for set B, and setting up the equation with $b$ and $e$ being the first and last terms of set $B$, 

 
 \[m(b+e) = m\]

 and so $b+e = 1$.

 Now we know, assume that both sequences are increasing sequences, for the sake of simplicity.  Based on the fact that set $A$ has half the number of elements as set $B$, and the difference between the greatest terms of the two two sequences is $99$ (forget about absolute value, it's insignificant here since we can just assume both sets end with positive last terms), you can set up an equation where $x$ is the last term of set A:

 
 \[2(x-(-x+4)+1) = 1+(x+99)-(-x-99+1)\]

 Note how i basically just counted the number of terms in each sequence here.  It's made a lot simpler because we just assumed that the first term is negative and last is positive for each set, it has absolutely no effect on the end result!  This is a great strategy that can help significantly simplify problems.  Also note how exactly i used the fact that the first and last terms of each sequence sum to $4$ and $1$ respectively (add $x$ and $(-x+4)$ to see what i mean).

 Solving this equation we find $x = 102$.  We know the first and last terms have to sum to $4$ so we find the first term of the sequence is $-98$.  Now, the solution is in clear sight, we just find the number of integers between $-98$ and $102$, inclusive, and it is $m = \boxed{201}$.

 Note how this method is not very algebra heavy.  It seems like a lot by the amount of text but really the first two steps are quite simple.

 



\textbf{Solution 4 (Sketchy Solution for Speedrunners)}

First, calculate the average of set $A$ and set $B$. It's obvious that they are $2$ and $1/2$ respectively. Let's look at both sets. Obviously, there is an odd number of integers in the set with $2$ being in the middle, which means that $m$ is an odd number and that the number of consecutive integers on each side of $2$ are equal. In set $B$, it is clear that it contains an even number of integers, but since the number in the middle is $1/2$, we know that the range of the consecutive numbers on both sides will be $(x$ to $0)$ and $(1$ to $-y)$. 

 Nothing seems useful right now, but let's try plugging an odd number, $3$, for $m$ in set $B$. We see that there are $6$ consecutive integers and $3$ on both sides of $1/2$. After plugging this into set $A$, we find that the set equals 
 \[{1,2,3}\]. From there, we find the absolute value of the difference of both of the greatest values, and get 0. 

 Let's try plugging in another odd number, $55$. We see that the resulting set of numbers is $(-54$ to $0)$, and $(1$ to  $55)$. We then plug this into set $A$, and find that the set of numbers is $(-25$ to $-29)$ which indeed results in the average being $2$. We then find the difference of the greatest values to be 26. 

 From here, we see a pattern that can be proven by more trial and error. When we make $m$ equal to $3$, then the difference is $0$ whearas when we make it $55$, then the difference is $26$. $55-3$ equals to $52$ and $26-0$ is just $0$. We then see that $m$ increases twice as fast as the difference. So when the difference is $99$, it increased $99$ from when it was $0$, which means that $m$ increased by $99*2$ which is $198$. We then add this to our initial $m$ of $3$, and get $\boxed{201}$ as our answer.

 



\textbf{Solution 5}

Let the first term of $A$ be $a$ and the first term of $B$ be $b$. There are $m$ elements in $A$ so $A$ is $a, a+1, a+2,...,a+m-1$. Adding these up, we get $\frac{2a+m-1}{2}\cdot m = 2m \implies 2a+m=5$. Set $B$ contains the numbers $b, b+1, b+2,...,b+2m-1$. Summing these up, we get $\frac{2b+2m-1}{2}\cdot 2m =m \implies 2b+2m=2$. The problem gives us that the absolute value of the difference of the largest terms in $A$ and $B$ is $99$. The largest term in $A$ is $a+m-1$ and the largest term in $B$ is $b+2m-1$ so $|b-a+m|=99$. From the first two equations we get, we can get that $2(b-a)+m=-3$. Now, we make a guess and assume that $b-a+m=99$ (if we get a negative value for $m$, we can try $b-a+m=-99$). From here we get that $b-a=-102$. Solving for $m$, we get that the answer is $\boxed{201}$

 -Heavytoothpaste

 



\textbf{Solution 6 (similar to all other solutions, gives intuition on )}

Let the first element of $\mathcal{A}$ be $p$ and the first element of $\mathcal{B}$ be $q$. 

 It is given $\mathcal{A}$ has $m$ members, and each member is represented as the sequence $\{p, p+1, p+2, \dots , p+m-1 \}$. It is given that $\mathcal{B}$ has $2m$ members, and each member is represented as the sequence $\{q, q+1, q+2, \dots, q+2m-1\}$.

 Now note that $\sum \mathcal{A} = 2m$ and $\sum \mathcal{B} = m$. It is seen that each of $\mathcal{A}$, $\mathcal{B}$ is an \textbf{arithmetic sequence}, and thus the sum $\mathcal{S}$ can be determined as

 
 \[\frac{2p+m-1}{2} \cdot m = 2m \implies 2p+m = 5\]
 \[\frac{2q+2m-1}{2} \cdot 2m = m \implies q+m = 1\]

 From the third statement, $|p - q - m| = 99$. Definition of absolute value is that $p-q-m = 99$ or $p-q-m = -99$.

 Applying the first equation gives $n = 100$ making $m$ negative, which is not allowed as counts cannot be negative. Otherwise, $n = -98$ and then $m = \boxed{201}$, a valid and satisfactory result.

 


	\item (\textit{AIME I 2004}) A convex polyhedron $P$ has 26 vertices, 60 edges, and 36 faces, 24 of which are triangular, and 12 of which are quadrilaterals. A space diagonal is a line segment connecting two non-adjacent vertices that do not belong to the same face. How many space diagonals does $P$ have?


 
 \textbf{Answer: }241



\textbf{Solution}

Every pair of \href{https://artofproblemsolving.com/wiki/index.php?title=Vertex}{ vertices} of the \href{https://artofproblemsolving.com/wiki/index.php?title=Polyhedron}{polyhedron} determines either an \href{https://artofproblemsolving.com/wiki/index.php?title=Edge}{edge}, a face \href{https://artofproblemsolving.com/wiki/index.php?title=Diagonal}{diagonal} or a space diagonal.  We have ${26 \choose 2} = \frac{26\cdot25}2 = 325$ total \href{https://artofproblemsolving.com/wiki/index.php?title=Line\_segment}{line segments} determined by the vertices.  Of these, $60$ are edges.  Each \href{https://artofproblemsolving.com/wiki/index.php?title=Triangle}{triangular} face has $0$ face diagonals and each \href{https://artofproblemsolving.com/wiki/index.php?title=Quadrilateral}{quadrilateral} face has $2$, so there are $2 \cdot 12 = 24$ face diagonals.  This leaves $325 - 60 - 24 = \boxed{241}$ segments to be the space diagonals.

 


	\item (\textit{AIME I 2004}) A square has sides of length 2. Set $S$ is the set of all line segments that have length 2 and whose endpoints are on adjacent sides of the square. The midpoints of the line segments in set $S$ enclose a region whose area to the nearest hundredth is $k.$ Find $100k.$

 \textbf{Answer: }086



\textbf{Solution}

Without loss of generality, let $(0,0)$, $(2,0)$, $(0,2)$, and $(2,2)$ be the \href{https://artofproblemsolving.com/wiki/index.php?title=Vertex}{ vertices} of the square. Suppose the endpoints of the segment lie on the two sides of the square determined by the vertex $(0,0)$.  Let the two endpoints of the segment have coordinates $(x,0)$ and $(0,y)$. Because the segment has length 2, $x^2+y^2=4$. Using the midpoint formula, we find that the midpoint of the segment has coordinates $\left(\frac{x}{2},\frac{y}{2}\right)$. Let $d$ be the distance from $(0,0)$ to $\left(\frac{x}{2},\frac{y}{2}\right)$. Using the \href{https://artofproblemsolving.com/wiki/index.php?title=Distance\_formula}{distance formula} we see that $d=\sqrt{\left(\frac{x}{2}\right)^2+\left(\frac{y}{2}\right)^2}= \sqrt{\frac{1}{4}\left(x^2+y^2\right)}=\sqrt{\frac{1}{4}(4)}=1$. Thus the midpoints lying on the sides determined by vertex $(0,0)$ form a quarter-\href{https://artofproblemsolving.com/wiki/index.php?title=Circle}{circle} with \href{https://artofproblemsolving.com/wiki/index.php?title=Radius}{radius} 1.  

 
                    \begin{figure}[H]
                        \centering
                        \includegraphics[width=0.5\textwidth]{Solution/AIME_I_2004_4_Solution_Figure_1.png}
                    \end{figure}
                    

 The set of all midpoints forms a quarter circle at each corner of the square.  The area enclosed by all of the midpoints is $4-4\cdot \left(\frac{\pi}{4}\right)=4-\pi \approx .86$ to the nearest hundredth. Thus $100\cdot k=\boxed{86}$




\textbf{Solution 2}


                    \begin{figure}[H]
                        \centering
                        \includegraphics[width=0.5\textwidth]{Solution/AIME_I_2004_4_Solution_Figure_2.png}
                    \end{figure}
                    

 If we imagine an arbitrary line with length $2$ connecting two sides of the square, we can draw the rectangle formed by drawing a perpendicular from where that line touches the square. 

 Drawing the other diagonal of the rectangle, it also has length two, and it bisects with the original line. Since their intersection is the midpoint of both lines, the distance from the corner to the midpoint is always $1$, which forms a circle with radius $1$ centered at the corner of the square. 

 The area of the shape then follows from simple calculations.

 


 



\textbf{Solution 3}

To imagine the area, think of a ladder with a length of $2$ sliding down a wall. It is known that as a ladder slides down a wall, its midpoint traces a quarter circle (if you don't believe me, try it with your pencil). There are $4$ quarter circles, so their area is one circle or $\pi$. Thus, they enclose the area of the square minus the area of the quarter circles, which is $4-\pi \approx 0.86$, so $100k = \boxed{086}$. ~Extremelysupercooldude

 



\textbf{Solution 4}

Consider the right triangle formed by the line segment and the vertex on the two sides it intersects. Since this right triangle has hypotenuse $2$, its circumradius is $1$. Since the circumcenter of a right triangle is the midpoint of the hypotenuse, the distance between the vertex of the square and the midpoint of the line segment is always $1$. Thus, $S$ forms $4$ quarter-circles of radius $1$ around the vertices of the square. Combined, these quarter-circles have total area of $\pi$, so the area enclosed by $S$ is $4-\pi \approx .86$, and the answer is $\boxed{86}$. ~A\_MatheMagician

 


	\item (\textit{AIME I 2004}) Alpha and Beta both took part in a two-day problem-solving competition. At the end of the second day, each had attempted questions worth a total of $500$ points. Alpha scored $160$ points out of $300$ points attempted on the first day, and scored $140$ points out of $200$ points attempted on the second day. Beta who did not attempt $300$ points on the first day, had a positive integer score on each of the two days, and Beta's daily success rate (points scored divided by points attempted) on each day was less than Alpha's on that day. Alpha's two-day success ratio was $\frac{300}{500} = \frac{3}{5}$. The largest possible two-day success ratio that Beta could achieve is $m/n,$ where $m$ and $n$ are relatively prime positive integers. What is $m+n$?


 
 \textbf{Answer: }849



\textbf{Solution}

Let $q$ be the number of questions Beta takes on day 1 and $a$ be the number he gets right. Let $b$ be the number he gets right on day 2.

 These inequalities follow:
 \[\frac{a}{q} < \frac{160}{300} = \frac{8}{15}\]
 \[\frac{b}{500-q} < \frac{140}{200} = \frac{7}{10}\]Solving for a and b and adding the two inequalities:
 \[a + b < \frac{8}{15}q + (350 - \frac{7}{10}q)\]
 \[a + b < 350 - \frac{1}{6}q\]From here, we see the largest possible value of $a+b$ is $349$.

 Checking our conditions, we know that $a$ must be positive so therefore $q$ must be positive. A quick check shows that for $2\le q \le 5$, $q$ follows all the conditions and results in $a+b=349$.

 This makes Beta's success ratio $\frac{349}{500}$. Thus, the answer is $m+n = 349 + 500 = \boxed{849}$.

 \textbf{Simple Solution} We see that we want Beta to have more points where there is a higher Alpha success rate (that way, the score is shifted more towards the higher Alpha score). With that in mind, $.7 > 8/15$. Thus, we might as well put as many points as possible into the second day.

 That said, we need Beta to have a positive integer score on either end per requirements. The best way is for Beta to go $\frac{1}{2}$ on the first day. That way, we "waste" the least points- both in the numerator and denominator.

 Thereafter, we see that, letting $x$ be the number of points Beta gained on the second day, $\frac{x}{498} < \frac{7}{10}$; thus $\max(x) = 348$.

 Aha. $348+1+500= \framebox{849}$.

 


	\item (\textit{AIME I 2004}) An integer is called snakelike if its decimal representation $a_1a_2a_3\cdots a_k$ satisfies $a_i<a_{i+1}$ if $i$ is odd and $a_i>a_{i+1}$ if $i$ is even. How many snakelike integers between 1000 and 9999 have four distinct digits?


 
 \textbf{Answer: }882



\textbf{Solution 1}

We divide the problem into two cases: one in which zero is not one of the digits and one in which it is. In the former case, suppose we pick digits $x_1,x_2,x_3,x_4$ such that $x_1<x_2<x_3<x_4$. There are five arrangements of these digits that satisfy the condition of being snakelike: $x_1x_3x_2x_4$, $x_1x_4x_2x_3$, $x_2x_3x_1x_4$, $x_2x_4x_1x_3$, $x_3x_4x_1x_2$. Thus there are $5\cdot {9\choose 4}=630$ snakelike numbers which do not contain the digit zero.

 In the second case we choose zero and three other digits such that $0<x_2<x_3<x_4$. There are three arrangements of these digits that satisfy the condition of being snakelike: $x_2x_30x_4$, $x_2x_40x_3$, $x_3x_40x_2$. Because we know that zero is a digit, there are $3\cdot{9\choose 3}=252$ snakelike numbers which contain the digit zero. Thus there are $630+252=\boxed{882}$ snakelike numbers.

 



\textbf{Solution 2}

Let's create the snakelike number from digits $a < b < c < d$, and, if we already picked the digits there are 5 ways to do so, as said in the first solution. And, let's just pick the digits from 0-9. This get's a total count of $5\cdot{10 \choose 4}$ But, this over-counts since it counts numbers like 0213. We can correct for this over-counting. Lock the first digit as 0 and permute 3 other chosen digits $a < b < c$. There are 2 ways to permute to make the number snakelike, b-a-c, or c-a-b. And, we pick a,b,c from 1 to 9, since 0 has already been chosen as one of the digits. So, the amount we have overcounted by is $2\cdot{9 \choose 3}$. Thus our answer is $5\cdot{10 \choose 4} - 2\cdot{9 \choose 3} = \boxed{882}$




\textbf{Solution 3}

We will first decide the order of the 4 digits, greatest to least. To do this, we will pretend that we have selected the digits 1,2,3,4, and we need to arrange them to create a snakelike number. By testing all permutations, there are only 5 ways to make a snakelike number: (1,3,2,4),(1,4,2,3),(2,3,1,4),(2,4,1,3),(3,4,1,2). 

 Now we select 4 digits to replace the 1,2,3,4. 

 In first 2 of cases: (1,3,2,4),(1,4,2,3), the leading digit is a 1, which means it is the smallest of our 4 digits. If we select a 0, the leading digit will be a zero, which is bad because all numbers between 1000 and 9999 have nonzero leading digits. So, we need to select our 4 digits only from the pool of 1-9. There are 9 choose 4 ways and there are 2 cases:

 $2 \cdot \dbinom{9}{4} = 252$

 Thus, there are 252 ways for those 2 cases. 

 For the next 3 cases, selecting a 0 is okay, so we can select from the pool of 0-9. There are 10 choose 4 ways to select our 4 digits and there are 3 cases:

 $3 \cdot \dbinom{10}{4} = 630$

 For those 3 cases there are 630 ways. 

 Thus, our answer is 630+252 = $\boxed{882}$.

 -Alexlikemath

 


	\item (\textit{AIME I 2004}) Let $C$ be the coefficient of $x^2$ in the expansion of the product $(1 - x)(1 + 2x)(1 - 3x)\cdots(1 + 14x)(1 - 15x).$ Find $|C|.$

 \textbf{Answer: }588



\textbf{Solution 1}

Let our \href{https://artofproblemsolving.com/wiki/index.php?title=Polynomial}{polynomial} be $P(x)$.

 It is clear that the coefficient of $x$ in $P(x)$ is $-1 + 2 - 3 + \ldots + 14 - 15 = -8$, so $P(x) = 1 -8x + Cx^2 + Q(x)$, where $Q(x)$ is some polynomial \href{https://artofproblemsolving.com/wiki/index.php?title=Divisibility}{ divisible} by $x^3$.

 Then $P(-x) = 1 + 8x + Cx^2 + Q(-x)$ and so $P(x)\cdot P(-x) = 1 + (2C - 64)x^2 + R(x)$, where $R(x)$ is some polynomial divisible by $x^3$.

 However, we also know $P(x)\cdot P(-x) = (1 - x)(1 + x)(1 +2x)(1 - 2x) \cdots (1 - 15x)(1 + 15x)$ $= (1 - x^2)(1 - 4x^2)\cdots(1 - 225x^2)$ $= 1 - (1 + 4 + \ldots + 225)x^2 + R(x)$.

 Equating coefficients, we have $2C - 64 = -(1 + 4 + \ldots + 225) = -1240$, so $-2C = 1176$ and $|C| = \boxed{588}$.

 



\textbf{Solution 2}

Let $S$ be the \href{https://artofproblemsolving.com/wiki/index.php?title=Set}{set} of integers $\{-1,2,-3,\ldots,14,-15\}$. The coefficient of $x^2$ in the expansion is equal to the sum of the product of each pair of distinct terms, or $C = \sum_{1 \le i \neq j}^{15} S_iS_j$. Also, we know that 
 \begin{align*}\left(\sum_{i=1}^{n} S_i\right)^2 &= \left(\sum_{i=1}^{n} S_i^2\right) + 2\left(\sum_{1 \le i \neq j}^{15} S_iS_j\right)\\ (-8)^2 &= \frac{15(15+1)(2\cdot 15+1)}{6} + 2C\end{align*}
where the left-hand sum can be computed from:

 $\sum_{i=1}^{15} S_i = S_{15} + \left(\sum_{i=1}^{7} S_{2i-1} + S_{2i}\right) = -15 + 7 = -8$

 and the right-hand sum comes from the formula for the sum of the first $n$ perfect squares. Therefore, $|C| = \left|\frac{64-1240}{2}\right| = \boxed{588}$.

 



\textbf{Solution 3 (Bash)}

Consider the set $[-1, 2,-3,4,-5,6,-7,8,-9,10,-11,12,-13,14,-15]$. Denote by $S$ all size 2 subsets of this set. Replace each element of $S$ by the product of the elements. Now, the quantity we seek is the sum of each element. Since consecutive elements add to $1$ or $-1$, we can simplify this to $588$




\textbf{Solution 4 (Easy)}

We know that this polynomial has roots $1, -\frac{1}{2}, \frac{1}{3},\ldots$ and the coefficient of $x^2$ will be the sum of the product taken by $13$. However, since this is closer to the constant side, we can create a new polynomial with the reciprocal roots which will make $C$ the coefficient of $x^{13}$ and thus the sum of the reciprocal roots taken by 2. We can calculate this with $1\cdot(8-1)-2\cdot(8-(-2))+\ldots+15(8-15)$. This will give us $2C = -1176$ giving us a final answer of $|C| = \fbox{588}$.

 


 



\textbf{Solution 4.5 (Easier Explanation)}

In this polynomial, we can see that by Vieta’s Formulae, the coefficient of $x^2$ is the sum of the products of every 13 of the roots. Let us denote the roots as $r_1, r_2, \dots, r_{15}$.

 


 Then, by Vieta’s, we have the sum:

 
 \[r_1 r_2 \dots r_{13} + \dots + r_3 r_4 \dots r_{15}\]

 


 Also, since we know that the product of the roots is $1$ (again by Vieta’s Formulae), we can rewrite an arbitrary term of this expression. For example, the term $r_1 r_2 \dots r_{13}$ is equal to:

 
 \[\frac{r_1 r_2 \dots r_{15}}{r_{14} r_{15}} = \frac{1}{r_{14} r_{15}}\]

 


 Now we are just finding the value of the expression, because this is just an expansion of:

 
 \[\frac{\left(\frac{1}{r_1} + \dots + \frac{1}{r_{15}}\right)^2 - \left(\frac{1}{r_1^2} + \dots + \frac{1}{r_{15}^2}\right)}{2}\]

 


 Which evaluates to:

 
 \[8 \cdot \frac{147}{2} = \boxed{588}\]




\textbf{Solution 4}

Let set $N$ be $\{-1, -3, \ldots -15\}$ and set $P$ be $\{2, 4, \ldots 14\}$. The sum of the negative $x^2$ coefficients is the sum of the products of the elements in all two element sets such that one element is from $N$ and the other is from $P$. Each summand is a term in the expansion of
 \[(-1 - 3 - \ldots - 15)(2 + 4 + \ldots + 14)\]which equals $-56 * 64 = -(60^2 - 4^2) = -3584$. The sum of the positive $x^2$ coefficients is the sum of the products of all two element sets such that the two elements are either both in $N$ or both in $P$. By counting, the sum is $2992$, so the sum of all $x^2$ coefficients is $-588$. Thus, the answer is $\boxed{588}$.

 


 



\textbf{Solution 5}

We can find out the coefficient of $x^2$ by multiplying every pair of two coefficients for $x$. This means that we multiply $-1$ by $2,-3,4,-5,6,-7,8,-9,10,-11,12,-13,14,-15$ and $2$ by $3,4,-5,6,-7,8,-9,10,-11,12,-13,14,-15$. and etc. This sum can be easily simplified and is equal to $(-1)(-7)+(-3)(-6)+(-5)(-5)+(-7)(-4)+(-9)(-3)+(-11)(-2)+(-13)(-1)+2(-9)+4(-10)+6(-11)+8(-12)+10(-13)+12(-14)+14(-15)$ or $588$.

 -David Camacho

 



\textbf{Solution 6}

This is just another way of summing the subsets of 2 from $[-1, 2, -3, 4, -5, 6, -7, 8, -9, 10, -11, 12, -13, 14, -15]$. Start from the right and multiply -15 to everything on its left. Use the distributive property and add all the 14 integers together to get 7. This gives us $-15 * 7$. Doing this for 14 gives us $14 * -7$, and for -13 we get $-13 * 6$. This pattern repeats where every two integers will multiple 7, 6,... to 0. Combining and simplifying the pattern give us this: $-(29 * 7 + 25 * 6 + 21 * 5 + 17 * 4 + 13 * 3 + 9 * 2 + 5*1)$. Calculate directly or notice that this sum is just $-\sum_{i=1}^{7} i(4i+1) = -4\sum_{i=1}^{7} i^2 -\sum_{i=1}^{7} i = -4\cdot\frac{(7)(8)(15)}{6} - \frac{(7)(8)}{2}$. The expression gives us -588, or $C = \boxed{588}$. This is a good solution because it guarantees we never add a product twice, and the pattern is simple to add by hand.

 -jackshi2006

 



\textbf{Solution 7}

We expand and obtain $\left(x-1\right)\left(2x-1\right)\left(3x-1\right)\cdots\left(15x-1\right) =  1307674368000x^{15} - 948550176000x^{14} - 689324826240x^{13} + 2733483288464x^{12} + 82808260416x^{11} - 23038684088x^{10} - 3811851848x^9 + 828730833x^8 + 81228128x^7 - 14661124x^6 - 853104x^5 + 132902x^4 + 4256x^3 - \boxed{588}x^2 - 8x + 1.$

 Do not do this in an actual competition.

 



\textbf{Solution 8 (Bash Yay!)}

Notice a pattern betweenthe x cooeficient and the x\textasciicircum{}2

 x {-1,1, -2, 2, -3, 3..}

 And then we can bash out the x\textasciicircum{}2 terms absolute value:

 x\textasciicircum{}2={0, -2, -5, -13,....} you get $\boxed{588}$




\textbf{Solution 9 (Recursion)}

Let $f(x)$ be the function that determines the previous coefficient of a polynomial in the form of the question for $x^2$. Define $g(x)$ analogously with coefficient of $x$. We see 

 
 \[f(1) = 1-x, f(1) = 0, g(1) = -1\] 
 \[f(2) = (1-x)(1+2x), f(2) = 2, g(2) = -3\] 
 \[f(n) \ : \ n \ge 3 = f(n-1) + g(n-1) \cdot (-1)^{n}nx\]. 
 \[g(n) \ : \ n \ge 3 = \text{sum of all} \ x \ \text{terms}\]. 

 From here, we can either find an explicit equation, or we can just bash (which will, unlike most other solutions, not take eternity).

 
 \[g(3) = -1+2-3 = -2, f(3) = -5\] 
 \[g(4) = 2, f(4) = -13\]
 \[g(5) = -3, f(5) = -23\]
 \[g(6) = 3, f(6) = -41\]
 \[g(7) = -4, f(7) = -62\]
 \[g(8) = 4, f(8) = -94\]
 \[g(9) = -5, f(9) = -130\]
 \[g(10) = 5, f(10) = -180\]
 \[g(11) = -6, f(11) = -235\]
 \[g(12) = 6, f(12) = -307\]
 \[g(13) = -7, f(13) = -385\]
 \[g(14) = 7, f(14) = -483\]
 \[g(15) = -8, f(15) = -588\].

 Note that $|f(15)| = |C|$. So the answer is $\boxed{588}$.

 


	\item (\textit{AIME I 2004}) Define a regular $n$-pointed star to be the union of $n$ line segments $P_1P_2, P_2P_3,\ldots, P_nP_1$ such that


 
            \begin{itemize}
        \item the points $P_1, P_2,\ldots, P_n$ are coplanar and no three of them are collinear,

        \item each of the $n$ line segments intersects at least one of the other line segments at a point other than an endpoint,

        \item all of the angles at $P_1, P_2,\ldots, P_n$ are congruent,

        \item all of the $n$ line segments $P_1P_2, P_2P_3,\ldots, P_nP_1$ are congruent, and

        \item the path $P_1P_2, P_2P_3,\ldots, P_nP_1$ turns counterclockwise at an angle of less than 180 degrees at each vertex.
            \end{itemize}
            There are no regular 3-pointed, 4-pointed, or 6-pointed stars. All regular 5-pointed stars are similar, but there are two non-similar regular 7-pointed stars. How many non-similar regular 1000-pointed stars are there?


 
 \textbf{Answer: }199



\textbf{Solution 1}

We use the \href{https://artofproblemsolving.com/wiki/index.php?title=Principle\_of\_Inclusion-Exclusion}{Principle of Inclusion-Exclusion} (PIE).

 If we join the adjacent vertices of the  regular $n$-star, we get a regular $n$-gon. We number the vertices of this $n$-gon in a counterclockwise direction: $0, 1, 2, 3, \ldots, n-1.$

 A regular $n$-star will be formed if we choose a vertex number $m$, where $0 \le m \le n-1$, and then form the line segments by joining the following pairs of vertex numbers:$(0 \mod{n}, m \mod{n}),$$(m \mod{n}, 2m \mod{n}),$$(2m \mod{n}, 3m \mod{n}),$$\cdots,$$((n-2)m \mod{n}, (n-1)m \mod{n}),$$((n-1)m \mod{n}, 0 \mod{n}).$

 If $\gcd(m,n) > 1$, then the star degenerates into a regular $\frac{n}{\gcd(m,n)}$-gon or a (2-vertex) line segment if$\frac{n}{\gcd(m,n)}= 2$. Therefore, we need to find all $m$ such that $\gcd(m,n) = 1$.

 Note that $n = 1000 = 2^{3}5^{3}.$

 Let $S = \{1,2,3,\ldots, 1000\}$, and $A_{i}= \{i \in S \mid i\, \textrm{ divides }\,1000\}$. The number of $m$'s that are not relatively prime to $1000$ is:$\mid A_{2}\cup A_{5}\mid = \mid A_{2}\mid+\mid A_{5}\mid-\mid A_{2}\cap A_{5}\mid$$= \left\lfloor \frac{1000}{2}\right\rfloor+\left\lfloor \frac{1000}{5}\right\rfloor-\left\lfloor \frac{1000}{2 \cdot 5}\right\rfloor$$= 500+200-100 = 600.$

 Vertex numbers $1$ and $n-1=999$ must be excluded as values for $m$ since otherwise a regular n-gon, instead of an n-star, is formed.

 The cases of a 1st line segment of (0, m) and (0, n-m) give the same star. Therefore we should halve the count to get non-similar stars.

 Therefore, the number of non-similar 1000 pointed stars is $\frac{1000-600-2}{2}= \boxed{199}.$

 Note that in general, the number of $n$-pointed stars is given by $\frac{\phi(n)}{2} - 1$ (dividing by $2$ to remove the reflectional symmetry, subtracting $1$ to get rid of the $1$-step case), where $\phi(n)$ is the \href{https://artofproblemsolving.com/wiki/index.php?title=Euler\%27s\_totient\_function}{Euler's totient function}. It is well-known that $\phi(n) = n\left(1-\frac{1}{p_1}\right)\left(1-\frac{1}{p_2}\right)\cdots \left(1-\frac{1}{p_n}\right)$, where $p_1,\,p_2,\ldots,\,p_n$ are the distinct prime factors of $n$. Thus $\phi(1000) = 1000\left(1 - \frac 12\right)\left(1 - \frac 15\right) = 400$, and the answer is $\frac{400}{2} - 1 = 199$.

 



\textbf{Solution 2}

Continue as before by cyclically constructing a star by taking every $k$-th point. We can think of this construction as the orbit of a point by a rotation $R^k\in D_{1000}$. Then, we want $|R^k|=1000$, meaning that $k\in (\mathbb Z/1000\mathbb Z)^\times$. Then, $k$ is coprime to $1000$, meaning $k$ may take on $\phi(1000)=400$ possible values. For any $R^k$, the dihedral symmetry $TR^k$ defines the figure-preserving bijection $R^k\mapsto R^{-k}$, meaning half of our possible values of $k$ are similar. Finally, the condition that each edge must intersect another rules out $k=1$, giving a total of $\frac{400}2-1=\boxed{199}$ pointed stars.

 



\textbf{Solution 3}

Consider the $n$-vertex regular polygon. This polygon has exactly $\frac{n(n-3)}{2}$ diagonals, excluding edges. Such a star is formed in the following manner.

 Choose a position point (a point that is fixed, say $P_0$). Consider a point, and draw the line from $P_0$ to that new point. Reflect this point over the line of symmetry down from $P_0$. Now, rotate this figure $n$ times to each vertex (obviously, not to the two points in which the segment was first drawn and $P_0$). The result formed should have exactly $n$ unique segments that form a unique star. The only way this can be possible is if no common point $P_{\not \varphi(n)}$ exists such that some mapping will map a diagonal back onto a previous diagonal. Therefore, all such stars must contain segments such that the number of external side lengths between the formed diagonal and its reflection of the polygon is relatively prime to $n$, or $\varphi(n)$ (the external side lengths are the lengths bounded by the two lines, shown in the figure below, an example of one of the two 7-sided stars.

 
                    \begin{figure}[H]
                        \centering
                        \includegraphics[width=0.5\textwidth]{Solution/AIME_I_2004_8_Solution_Figure_1.png}
                    \end{figure}
                    

 Now, every star can be reflected over this line of symmetry, and we end up with a star that is also apart of $\varphi(n)$. Therefore, we must divide by 2. (In other terms, we drew a line from $P_0$ to another point $P_{a}$, and we do the same for the point $P'_{a}$ reflected over the line of symmetry through $P_0$. That result is the same star overcounted double). So, we have $\frac{\varphi(n)}{2}$. Are we done? 

 No! Notice when we draw the star, we also have that $n-1$ is relatively prime to $n$ for all $n$. However, this star is the polygon itself, and that defies our definition of a star, because each segment must have one intersection such that the exterior angle is less than 180. In the case of a polygon, this is never the case, and we already have one copy remaining when we divide by 2.

 One might consider the "zero pointed star". However, because 0 is not defined to be relatively prime to any number (you cannot divide by 0), we neglect the case entirely.

 Therefore, the total number of stars that can be formed given an $n$-gon is $\frac{\varphi(n)}{2} - 1$. Plugging in 1000, we have $\boxed{199}$.

 


	\item (\textit{AIME I 2004}) Let $ABC$ be a triangle with sides 3, 4, and 5, and $DEFG$ be a 6-by-7 rectangle. A segment is drawn to divide triangle $ABC$ into a triangle $U_1$ and a trapezoid $V_1$ and another segment is drawn to divide rectangle $DEFG$ into a triangle $U_2$ and a trapezoid $V_2$ such that $U_1$ is similar to $U_2$ and $V_1$ is similar to $V_2.$ The minimum value of the area of $U_1$ can be written in the form $m/n,$ where $m$ and $n$ are relatively prime positive integers. Find $m+n.$

 \textbf{Answer: }035



\textbf{Solution}

We let $AB=3, AC=4, DE=6, DG=7$ for the purpose of labeling. Clearly, the dividing segment in $DEFG$ must go through one of its vertices, \href{https://artofproblemsolving.com/wiki/index.php?title=Without\_loss\_of\_generality}{without loss of generality} $D$. The other endpoint ($D'$) of the segment can either lie on $\overline{EF}$ or $\overline{FG}$. $V_2$ is a trapezoid with a right angle then, from which it follows that $V_1$ contains one of the right angles of $\triangle ABC$, and so $U_1$ is similar to $ABC$. Thus $U_1$, and hence $U_2$, are $3-4-5\,\triangle$s.  

 Suppose we find the ratio $r$ of the smaller base to the larger base for $V_2$, which consequently is the same ratio for $V_1$. By similar triangles, it follows that $U_1 \sim \triangle ABC$ by the same ratio $r$, and since the ratio of the areas of two similar figures is equal to the square of the ratio of their corresponding lengths, it follows that $[U_1] = r^2 \cdot [ABC] = 6r^2$.

 <table><tbody><tr><td>
                    \begin{figure}[H]
                        \centering
                        \includegraphics[width=0.5\textwidth]{Solution/AIME_I_2004_9_Solution_Figure_1.png}
                    \end{figure}
                    
                    \begin{figure}[H]
                        \centering
                        \includegraphics[width=0.5\textwidth]{Solution/AIME_I_2004_9_Solution_Figure_2.png}
                    \end{figure}
                    
</td></tr><tr><td>
                    \begin{figure}[H]
                        \centering
                        \includegraphics[width=0.5\textwidth]{Solution/AIME_I_2004_9_Solution_Figure_3.png}
                    \end{figure}
                    
                    \begin{figure}[H]
                        \centering
                        \includegraphics[width=0.5\textwidth]{Solution/AIME_I_2004_9_Solution_Figure_4.png}
                    \end{figure}
                    
</td></tr></tbody></table>

            \begin{itemize}
        \item If $D'$ lies on $\overline{EF}$, then $ED' = \frac34 \cdot 6 = \frac{9}{2},\, 8$ from maintaining the $3-4-5$ triangular ratio; the latter ($\frac43 \cdot 6 = 8$) can be discarded as extraneous since we need $ED'<EF$. Therefore, $D'F = \frac{5}{2}$, and the ratio $r = \frac{D'F}{DG} = \frac{5}{14}$. The area of $[U_1] = 6\left(\frac{5}{14}\right)^2$ in this case.
            \end{itemize}
            
            \begin{itemize}
        \item If $D'$ lies on $\overline{FG}$, then $GD' = \frac{21}{4},\, \frac{28}{3}$; the latter can be discarded as extraneous since we need $GD'<GF$. Therefore, $D'F = \frac{3}{4}$, and the ratio $r = \frac{D'F}{DE} = \frac{1}{8}$. The area of $[U_1] = 6\left(\frac{1}{8}\right)^2$ in this case.
            \end{itemize}
            Of the two cases, the second is smaller; the answer is $\frac{3}{32}$, and $m+n = \boxed{035}$.

 


	\item (\textit{AIME I 2004}) A circle of radius 1 is randomly placed in a 15-by-36 rectangle $ABCD$ so that the circle lies completely within the rectangle. Given that the probability that the circle will not touch diagonal $AC$ is $m/n,$ where $m$ and $n$ are relatively prime positive integers, find $m + n.$

 \textbf{Answer: }817



\textbf{Solution}


                    \begin{figure}[H]
                        \centering
                        \includegraphics[width=0.5\textwidth]{Solution/AIME_I_2004_10_Solution_Figure_1.png}
                    \end{figure}
                    




\textbf{Solution 1}

The location of the center of the circle must be in the $34 \times 13$ rectangle that is one unit away from the sides of rectangle $ABCD$. We want to find the area of the \href{https://artofproblemsolving.com/wiki/index.php?title=Right\_triangle}{right triangle} with \href{https://artofproblemsolving.com/wiki/index.php?title=Hypotenuse}{hypotenuse} one unit away from $\overline{AC}$. Let this triangle be $A'B'C'$.

 Notice that $ABC$ and $A'B'C'$ share the same \href{https://artofproblemsolving.com/wiki/index.php?title=Incenter}{incenter}; this follows because the corresponding sides are parallel, and so the perpendicular \href{https://artofproblemsolving.com/wiki/index.php?title=Inradius}{inradii} are concurrent, except that the inradii of $\triangle ABC$ extend one unit farther than those of $\triangle A'B'C'$. From $A = rs$, we note that $r_{ABC} = \frac{[ABC]}{s_{ABC}} = \frac{15 \cdot 36 /2}{(15+36+39)/2} = 6$. Thus $r_{A'B'C'} = r_{ABC} - 1 = 5$, and since the ratio of the areas of two similar figures is equal to the square of the ratio of their corresponding lengths, $[A'B'C'] = [ABC] \cdot \left(\frac{r_{A'B'C'}}{r_{ABC}}\right)^2 = \frac{15 \times 36}{2} \cdot \frac{25}{36} = \frac{375}{2}$. 

 The probability is $\frac{2[A'B'C']}{34 \times 13} = \frac{375}{442}$, and $m + n = \boxed{817}$.  

 



\textbf{Solution 2}


                    \begin{figure}[H]
                        \centering
                        \includegraphics[width=0.5\textwidth]{Solution/AIME_I_2004_10_Solution_Figure_1.png}
                    \end{figure}
                    

 Let the bisector of $\angle CAD$ be $AE$, with $E$ on $CD$. By the angle bisector theorem, $DE = 36/5$. Since $\triangle AOR \sim \triangle AED$ ($O$ is the center of the circle), we find that $AR = 5$ since $OR = 1$. Also $AT = 35$ so $RT = OQ = 30$.

 We can apply the same principle again to find that $PT = 27/2$, and since $QT = 1$, we find that $PQ = 27/2 - 1 = 25/2$. The locus of all possible centers of the circle on this "half" of the rectangle is triangle $\triangle OPQ$. There exists another congruent triangle that is symmetric over $AC$ that has the same area as triangle $\triangle OPQ$. $\triangle APQ$ has area $\frac {1}{2}\cdot OP \cdot PQ = \frac {1}{2}\cdot 30\cdot \frac {25}{2}$, since $\angle OQP$ is right. Thus the total area that works is $30\cdot \frac {25}{2} = 375$, and the area of the locus of all centers of any circle with radius 1 is $34\cdot 13 = 442$. Hence, the desired probability is $\frac {375}{442}$, and our answer is $\boxed {817}$.

 



\textbf{Solution 3}


                    \begin{figure}[H]
                        \centering
                        \includegraphics[width=0.5\textwidth]{Solution/AIME_I_2004_10_Solution_Figure_3.png}
                    \end{figure}
                    

 Again, the location of the center of the circle must be in the $34 \times 13$ rectangle that is one unit away from the sides of rectangle $ABCD$. We want to find the area of the \href{https://artofproblemsolving.com/wiki/index.php?title=Right\_triangle}{right triangle} with \href{https://artofproblemsolving.com/wiki/index.php?title=Hypotenuse}{hypotenuse} one unit away from $\overline{AC}$.

 Let $A$ be at the origin, $B (36,0)$, $C (36,15)$, $D (0,15)$. The slope of $\overline{AC}$ is $\frac{15}{36} = \frac{5}{12}$. Let $\triangle A'B'C'$ be the right triangle with sides one unit inside $\triangle ABC$. Since $\overline{AC} || \overline{A'C'}$, they have the same slope, and the equation of $A'C'$ is $y = \frac{5}{12}x + c$. Manipulating, $5x - 12y + 12c = 0$. We need to find the value of $c$, which can be determined since $\overline{AC}$ is one unit away from $\overline{A'C'}$. Since the diagonal contains the origin, we can use the distance from a point to the line formula at the origin:

 
 \[\left|\frac{Ax + By + C}{\sqrt{A^2+B^2}}\right| = 1 \Longrightarrow \left|\frac{(5)(0) + (-12)(0) + 12c}{\sqrt{5^2 + (-12)^2}}\right| = 1\]
 \[c = \pm \frac{13}{12}\]

 The two values of $c$ correspond to the triangle on top and below the diagonal. We are considering $A'B'C'$ which is below, so $c = -\frac{13}{12}$. Then the equation of $\overline{A'C'}$ is $y = \frac{5}{12}x - \frac{13}{12}$. Solving for its intersections with the lines $y = 1, x = 35$ (boundaries of the internal rectangle), we find the coordinates of $A'B'C'$ are at $A'\ (5,1)\ B'\ (35,1)\ C'\ (35,\frac{27}{2})$. The area is $\frac{1}{2}bh = \frac{1}{2}(35-5)\left(\frac{27}{2} - 1\right) = \frac{375}{2}$.

 Finally, the probability is $\frac{2\cdot \mathrm{area\ of\ triangle}}{34 \times 13} = \frac{375}{442}$, and $m + n = \boxed{817}$.

 For this solution, if you didn't know the formula for the distance from a point to a line, you can use similar triangles to get the ratio:

 $\frac{1}{36}=\frac{c}{39}$

 This yields $c=\frac{13}{12}$.

 


	\item (\textit{AIME I 2004}) A solid in the shape of a right circular cone is 4 inches tall and its base has a 3-inch radius. The entire surface of the cone, including its base, is painted. A plane parallel to the base of the cone divides the cone into two solids, a smaller cone-shaped solid $C$ and a frustum-shaped solid $F,$ in such a way that the ratio between the areas of the painted surfaces of $C$ and $F$ and the ratio between the volumes of $C$ and $F$ are both equal to $k.$ Given that $k=m/n,$ where $m$ and $n$ are relatively prime positive integers, find $m+n.$

 \textbf{Answer: }512



\textbf{Solution 1}

Our original solid has volume equal to $V = \frac13 \pi r^2 h = \frac13 \pi 3^2\cdot 4 = 12 \pi$ and has \href{https://artofproblemsolving.com/wiki/index.php?title=Surface\_area}{surface area} $A = \pi r^2 + \pi r \ell$, where $\ell$ is the \href{https://artofproblemsolving.com/wiki/index.php?title=Slant\_height}{slant height} of the cone.  Using the \href{https://artofproblemsolving.com/wiki/index.php?title=Pythagorean\_Theorem}{Pythagorean Theorem}, we get $\ell = 5$ and $A = 24\pi$.

 Let $x$ denote the \href{https://artofproblemsolving.com/wiki/index.php?title=Radius}{radius} of the small cone. Let $A_c$ and 

 


	\item (\textit{AIME I 2004}) Let $S$ be the set of ordered pairs $(x, y)$ such that $0 < x \le 1, 0<y\le 1,$ and $\left \lfloor{\log_2{\left(\frac 1x\right)}}\right \rfloor$ and $\left \lfloor{\log_5{\left(\frac 1y\right)}}\right \rfloor$ are both even. Given that the area of the graph of $S$ is $m/n,$ where $m$ and $n$ are relatively prime positive integers, find $m+n.$ The notation $\left \lfloor{z}\right \rfloor$ denotes the greatest integer that is less than or equal to $z.$

 \textbf{Answer: }014



\textbf{Solution}

$\left\lfloor\log_2\left(\frac{1}{x}\right)\right\rfloor$ is even when

 
 \[x \in \left(\frac{1}{2},1\right] \cup \left(\frac{1}{8},\frac{1}{4}\right] \cup \left(\frac{1}{32},\frac{1}{16}\right] \cup \cdots\]

 Likewise:$\left\lfloor\log_5\left(\frac{1}{y}\right)\right\rfloor$ is even when

 
 \[y \in \left(\frac{1}{5},1\right] \cup \left(\frac{1}{125},\frac{1}{25}\right] \cup \left(\frac{1}{3125},\frac{1}{625}\right] \cup \cdots\]

 Graphing this yields a series of \href{https://artofproblemsolving.com/wiki/index.php?title=Rectangle}{rectangles} which become smaller as you move toward the \href{https://artofproblemsolving.com/wiki/index.php?title=Origin}{origin}. The $x$ interval of each box is given by the \href{https://artofproblemsolving.com/wiki/index.php?title=Geometric\_sequence}{geometric sequence} $\frac{1}{2} , \frac{1}{8}, \frac{1}{32}, \cdots$, and the $y$ interval is given by $\frac{4}{5} , \frac{4}{125}, \frac{4}{3125}, \cdots$

 Each box is the product of one term of each sequence. The sum of these boxes is simply the product of the sum of each sequence or:

 
 \[\left(\frac{1}{2} + \frac{1}{8} + \frac{1}{32} \ldots \right)\left(\frac{4}{5} + \frac{4}{125} + \frac{4}{3125} \ldots \right)=\left(\frac{\frac{1}{2}}{1 - \frac{1}{4}}\right)\left(\frac{\frac{4}{5}}{1-\frac{1}{25}}\right)= \frac{2}{3} \cdot \frac{5}{6} = \frac{5}{9},\] and the answer is $m+n = 5 + 9 = \boxed{014}$.

 


	\item (\textit{AIME I 2004}) The polynomial $P(x)=(1+x+x^2+\cdots+x^{17})^2-x^{17}$ has 34 complex roots of the form $z_k = r_k[\cos(2\pi a_k)+i\sin(2\pi a_k)], k=1, 2, 3,\ldots, 34,$ with $0 < a_1 \le a_2 \le a_3 \le \cdots \le a_{34} < 1$ and $r_k>0.$ Given that $a_1 + a_2 + a_3 + a_4 + a_5 = m/n,$ where $m$ and $n$ are relatively prime positive integers, find $m+n.$

 \textbf{Answer: }482



\textbf{Solution}

We see that the expression for the \href{https://artofproblemsolving.com/wiki/index.php?title=Polynomial}{polynomial} $P$ is very difficult to work with directly, but there is one obvious transformation to make: sum the \href{https://artofproblemsolving.com/wiki/index.php?title=Geometric\_series}{geometric series}:

 
 \begin{align*} P(x) &= \left(\frac{x^{18} - 1}{x - 1}\right)^2 - x^{17} = \frac{x^{36} - 2x^{18} + 1}{x^2 - 2x + 1} - x^{17}\\ &= \frac{x^{36} - x^{19} - x^{17} + 1}{(x - 1)^2} = \frac{(x^{19} - 1)(x^{17} - 1)}{(x - 1)^2} \end{align*}


 This \href{https://artofproblemsolving.com/wiki/index.php?title=Expression}{expression} has \href{https://artofproblemsolving.com/wiki/index.php?title=Root}{roots} at every $17$th root and $19$th \href{https://artofproblemsolving.com/wiki/index.php?title=Roots\_of\_unity}{roots of unity}, other than $1$.  Since $17$ and $19$ are \href{https://artofproblemsolving.com/wiki/index.php?title=Relatively\_prime}{relatively prime}, this means there are no duplicate roots.  Thus, $a_1, a_2, a_3, a_4$ and $a_5$ are the five smallest fractions of the form $\frac m{19}$ or $\frac n {17}$ for $m, n > 0$.  

 $\frac 3 {17}$ and $\frac 4{19}$ can both be seen to be larger than any of $\frac1{19}, \frac2{19}, \frac3{19}, \frac 1{17}, \frac2{17}$, so these latter five are the numbers we want to add.

 $\frac1{19}+ \frac2{19}+ \frac3{19}+ \frac 1{17}+ \frac2{17}= \frac6{19} + \frac 3{17} = \frac{6\cdot17 + 3\cdot19}{17\cdot19} = \frac{159}{323}$ and so the answer is $159 + 323 = \boxed{482}$.

 



\textbf{Solution 2 (similar to 1)}


 \begin{align*} P(x) &= \left(\frac{x^{18} - 1}{x - 1}\right)^2 - x^{17} = \frac{x^{36} - 2x^{18} + 1}{x^2 - 2x + 1} - x^{17}\\ &= \frac{x^{36} - x^{19} - x^{17} + 1}{(x - 1)^2} = \frac{(x^{19} - 1)(x^{17} - 1)}{(x - 1)^2} \end{align*}


 Hence,

 
 \[(x-1)^2P(x) = (x^{19}-1)(x^{17}-1)\].

 Note $x \neq 1$ hence we just require some $x$ so that one of $x^{19}-1, x^{17}-1 =0$. That is, all 19th roots and 17th roots of unity.  This is because if $x=1$ we divide not only by 0 at the beginning, but also $x-1 = 0$ hence we have an indeterminable form. Then, each root is of the form:

 
 \[e^{\frac{2i\pi}{17}}, e^{\frac{2i\pi}{19}} \implies \operatorname{cis}(\frac{2k\pi}{17}), \operatorname{cis}(\frac{2k\pi}{19})\].

 Only consider the values of $\frac{k}{19}, \frac{k}{17}$. In order, the smallest of them all are, respectively, 
 \[\alpha_1, \alpha_2, \alpha_3, \alpha_4, \alpha_5 = \frac{1}{19}, \frac{1}{17}, \frac{2}{19}, \frac{2}{17}, \frac{3}{19}\].
 \[\sum_{i=1}^5 \alpha_i = \frac{159}{323} \implies \boxed{482} \quad \square\].

 


	\item (\textit{AIME I 2004}) A unicorn is tethered by a 20-foot silver rope to the base of a magician's cylindrical tower whose radius is 8 feet. The rope is attached to the tower at ground level and to the unicorn at a height of 4 feet. The unicorn has pulled the rope taut, the end of the rope is 4 feet from the nearest point on the tower, and the length of the rope that is touching the tower is $\frac{a-\sqrt{b}}c$ feet, where $a, b,$ and $c$ are positive integers, and $c$ is prime. Find $a+b+c.$

 \textbf{Answer: }813



\textbf{Solution}


                    \begin{figure}[H]
                        \centering
                        \includegraphics[width=0.5\textwidth]{Solution/AIME_I_2004_14_Solution_Figure_1.png}
                    \end{figure}
                    

 Looking from an overhead view, call the \href{https://artofproblemsolving.com/wiki/index.php?title=Center}{center} of the \href{https://artofproblemsolving.com/wiki/index.php?title=Circle}{circle} $O$, the tether point to the unicorn $A$ and the last point where the rope touches the tower $B$.  $\triangle OAB$ is a \href{https://artofproblemsolving.com/wiki/index.php?title=Right\_triangle}{right triangle} because $OB$ is a radius and $BA$ is a \href{https://artofproblemsolving.com/wiki/index.php?title=Tangent\_line}{tangent line} at point $B$.  We use the \href{https://artofproblemsolving.com/wiki/index.php?title=Pythagorean\_Theorem}{Pythagorean Theorem} to find the horizontal component of $AB$ has length $4\sqrt{5}$.  

 
                    \begin{figure}[H]
                        \centering
                        \includegraphics[width=0.5\textwidth]{Solution/AIME_I_2004_14_Solution_Figure_2.png}
                    \end{figure}
                    

 Now look at a side view and "unroll" the cylinder to be a flat surface. Let $C$ be the bottom tether of the rope, let $D$ be the point on the ground below $A$, and let $E$ be the point directly below $B$.  \href{https://artofproblemsolving.com/wiki/index.php?title=Triangle}{Triangles} $\triangle CDA$ and $\triangle CEB$ are \href{https://artofproblemsolving.com/wiki/index.php?title=Similar}{similar} \href{https://artofproblemsolving.com/wiki/index.php?title=Right\_triangle}{right triangles}.  By the Pythagorean Theorem $CD=8\cdot\sqrt{6}$.  

 Let $x$ be the length of $CB$. 
 \[\frac{CA}{CD}=\frac{CB}{CE}\implies \frac{20}{8\sqrt{6}}=\frac{x}{8\sqrt{6}-4\sqrt{5}}\implies x=\frac{60-\sqrt{750}}{3}\] 

 Therefore $a=60, b=750, c=3, a+b+c=\boxed{813}$.

 



\textbf{Solution 2}

Note that by Power of a Point, the point the unicorn is at has power $4 \cdot 20 = 80$ which implies that the tangent from that point to the tower is of length $\sqrt{80}=4\sqrt{5},$ however this is length of the rope projected into 2-D. If we let $\theta$ be the angle between the horizontal and the rope, we have that $\cos\theta=\frac{1}{5}$ which implies that $\sin\theta=\frac{2\sqrt{6}}{5}.$ Note that the portion of rope not on the tower is $4\sqrt{5} \cdot \frac{5}{2\sqrt{6}}= \frac{5\sqrt{30}}{3},$ the requested length of rope is $20-\frac{5\sqrt{30}}{3}=\frac{60-\sqrt{750}}{3}$ thus the requested sum is $\boxed{813}.$



	\item (\textit{AIME I 2004}) For all positive integers $x$, let

 \[f(x)=\begin{cases}1 &\mbox{if }x = 1\\ \frac x{10} &\mbox{if }x\mbox{ is divisible by 10}\\ x+1 &\mbox{otherwise}\end{cases}\]
and define a sequence as follows: $x_1=x$ and $x_{n+1}=f(x_n)$ for all positive integers $n$. Let $d(x)$ be the smallest $n$ such that $x_n=1$. (For example, $d(100)=3$ and $d(87)=7$.) Let $m$ be the number of positive integers $x$ such that $d(x)=20$. Find the sum of the distinct prime factors of $m$.


 
 \textbf{Answer: }511



\textbf{Solution 1}

We backcount the number of ways. Namely, we start at $x_{20} = 1$, which can only be reached if $x_{19} = 10$, and then we perform $18$ operations that either consist of $A: (-1)$ or $B: (\times 10)$. We represent these operations in a string format, starting with the operation that sends $f(x_{18}) = x_{19}$ and so forth downwards. There are $2^9$ ways to pick the first $9$ operations; however, not all $9$ of them may be $A: (-1)$ otherwise we return back to $x_{10} = 1$, contradicting our assumption that $20$ was the smallest value of $n$. Using \href{https://artofproblemsolving.com/wiki/index.php?title=Complementary\_counting}{complementary counting}, we see that there are only $2^9 - 1$ ways.

 Since we performed the operation $B: (\times 10)$ at least once in the first $9$ operations, it follows that $x_{10} \ge 20$, so that we no longer have to worry about reaching $1$ again. 

 However, we must also account for a sequence of $10$ or more $A: (-1)$s in a row, because that implies that at least one of those numbers was divisible by $10$, yet the $\frac{x}{10}$ was never used, contradiction. We must use complementary counting again by determining the number of strings of $A,B$s of length $18$ such that there are $10$ $A$s in a row. The first ten are not included since we already accounted for that scenario above, so our string of $10$ $A$s must be preceded by a $B$. Afterwards, if there was an $A$ before the string of $A$'s the string would include that letter. There are no other restrictions on the remaining seven characters. Letting $\square$ to denote either of the functions, and $^{[k]}$ to indicate that the character appears $k$ times in a row, our bad strings can take the forms:
 \begin{align*} &\underbrace{BA^{[10]}}\square \square \square \square \square \square \square\\ &\square \underbrace{BA^{[10]}}\square \square \square \square \square \square \\ &\qquad \quad \cdots \quad \cdots \\ &\square \square \square \square \square \square \underbrace{BA^{[10]}} \square \\ &\square \square \square \square \square \square \square \underbrace{BA^{[10]}} \\ \end{align*}


 There are $2^7$ ways to select the operations for the $7$ $\square$s, and $8$ places to place our $BA^{[10]}$ block. Thus, our answer is $2^9(2^9-1)-8\cdot 2^7 = 2^{18} - 2^9 - 8 \cdot 2^7 = 2^9 \times 509$, and the answer is $\boxed{511}$.

 \textbf{Note:} This solution is quick and most similar to the official solution; however, neither this nor the official solution prove that the final results of these inverted operations are all distinct. A more sophisticated argument, such as the one below, is needed to do so.

 



\textbf{Solution 2}

We approach the problem by \href{https://artofproblemsolving.com/wiki/index.php?title=Recursion}{recursion}. We \href{https://artofproblemsolving.com/wiki/index.php?title=Partition}{partition} the positive integers into the sets
 \[\mathcal{A}_{n,k}=\{x\in\mathbb{Z}^+\,:\, d(x)=n\text{ and } x\equiv k\pmod{10}\}.\]First, we note that $\mathcal{A}_{1,1}=\{1\}$, so by the \href{https://artofproblemsolving.com/wiki/index.php?title=Disjoint}{disjointness} of the $\mathcal{A}_{n,k}$'s, we know that $1$ is not in any of the other sets. Also, we note that $\mathcal{A}_{1,k}=\emptyset$ for $k=0,2,3,4,5,6,7,8,9$.

 We claim that if $2\le k\le 9$ and $n\ge2$, then $|\mathcal{A}_{n,k}|=|\mathcal{A}_{n-1,k+1}|$. To prove this, we show that $f$ (the given function) maps $\mathcal{A}_{n,k}$ bijectively onto $\mathcal{A}_{n-1,k+1}$. If $x\in \mathcal{A}_{n,k}$, then $x\equiv k\pmod{10}$, and $x\ne 1$, so $f(x)=x+1\equiv k+1\pmod{10}$. Also, $f^{(n)}(x)=1$, where $n$ is the smallest positive integer for which this is true. Therefore, $f^{{n-1}}(f(x))=1$, where $n-1$ is the smallest integer for which this is true. Thus $f(x)\in\mathcal{A}_{n-1,k+1}$. Also, since $f(x)=x+1$ on this set, we see that $f(x)=f(y)$ implies that $x=y$. Hence $f$ is an \href{https://artofproblemsolving.com/wiki/index.php?title=Injection}{injection}. If $y\in\mathcal{A}_{n-1,k+1}$, then $y-1\equiv k\pmod{10}$, where $2\le k\le 9$, so we know that $f(y-1)=y$, and $y-1\in \mathcal{A}_{n,k}$. Therefore, $f$ is a \href{https://artofproblemsolving.com/wiki/index.php?title=Surjection}{surjection}, so it must be a \href{https://artofproblemsolving.com/wiki/index.php?title=Bijection}{bijection}. Therefore,  if $2\le k\le 9$ and $n\ge2$, then $|\mathcal{A}_{n,k}|=|\mathcal{A}_{n-1,k+1}|$.

 We also claim that if $k=1$, $n\ge 2$ and $n\ne 11$, then  $|\mathcal{A}_{n,k}|=|\mathcal{A}_{n-1,k+1}|$. The proof is the same as in the previous paragraph, though some additional clarification is needed to prove that $f$ is a surjection. If $y\in\mathcal{A}_{n-1,k+1}$, or rather $y-1\equiv k\equiv 1\pmod{10}$, then if $y=2$, we see that $y-1=1$, and then $f(y-1)=1$, not $y$ as in the prior argument. However, this does not happen if $n\ne 11$. It is easy to check that $2\in\mathcal{A}_{10,2}$. Therefore, the only time that the above argument could fail is if $n-1=10$ and $k+1=2$. But in every other case, $|\mathcal{A}_{n,k}|=|\mathcal{A}_{n-1,k+1}|$.

 Next, we claim that if $n\ge 3$, then
 \[|\mathcal{A}_{n,0}|=\Big|\bigcup_{k=0}^9\mathcal{A}_{n-1,k}\Big|\]If $x\in\mathcal{A}_{n,0}$, then $f(x)=\frac{x}{10}$, which is clearly an injective map.  Also, $f^{(n)}(x)=1$, where $n$ is the smallest positive integer for which this is true. Therefore, $f^{{n-1}}(f(x))=1$, where $n-1$ is the smallest integer for which this is true. Thus $f(x)\in\mathcal{A}_{n-1,k}$ for some $k$. Conversely, if $y\in \mathcal{A}_{n-1,k}$, then $f(10y)=y$, so $d(10y)=n$, so $10y\in\mathcal{A}_{n,0}$.

 Based on these bijections, if we let $A_{n,k}=|\mathcal{A}_{n,k}|$, then
 \begin{align*} A_{n,0}&=\sum_{k=0}^{9} A_{n-1,k}\\ A_{n,1}&=A_{n-1,2}\qquad\text{if }n\ne 11\\ A_{n,2}&=A_{n-1,3}\\ A_{n,3}&=A_{n-1,4}\\ &\vdots\\ A_{n,8}&=A_{n-1,9}\\ A_{n,9}&=A_{n-1,0}. \end{align*}
Let $S_n=\sum_{k=0}^9 A_{n,k}$. Then by adding the above equations (valid if $n\ne 11$), we find that
 \[S_n=2S_{n-1}-A_{n-1,1}.\]Now by using the above relations repeatedly, we find
 \[A_{n-1,1}=A_{n-2,2}=A_{n-3,3}=\cdots=A_{n-9,9}=A_{n-10,0}=S_{n-11}.\]The first relation will only be valid if $n\ne 12$. Therefore, $S_n=2S_{n-1}-S_{n-11}$ for $n\ge 13$. For smaller values, it is easy to use the relations to compute that the terms are powers of $2$, although we note that $A_{11,1}=0$ because of the failure of the above argument.
 \[\begin{array}{c|cccccccccc|c} n&A_{n,1}&A_{n,2}&A_{n,3}&A_{n,4}&A_{n,5}&A_{n,6}&A_{n,7}&A_{n,8}&A_{n,9}&A_{n,0}&S_{n}\\\hline 1&1&0&0&0&0&0&0&0&0&0&1\\ 2&0&0&0&0&0&0&0&0&0&1&1\\ 3&0&0&0&0&0&0&0&0&1&1&2\\ 4&0&0&0&0&0&0&0&1&1&2&4\\ 5&0&0&0&0&0&0&1&1&2&4&8\\ 6&0&0&0&0&0&1&1&2&4&8&16\\ 7&0&0&0&0&1&1&2&4&8&16&32\\ 8&0&0&0&1&1&2&4&8&16&32&64\\ 9&0&0&1&1&2&4&8&16&32&64&128\\ 10&0&1&1&2&4&8&16&32&64&128&256\\ 11&\boxed{0}&1&2&4&8&16&32&64&128&256&511\\ 12&1&2&4&8&16&32&64&128&256&511&1022\\ \end{array}\]After this, we can simply use the recurrence relation for $S_n$, finding
 \[\begin{array}{c|l} n&S_n\\\hline 1&1\\ 2&1\\ 3&2^1\\ 4&2^2\\ 5&2^3\\ 6&2^4\\ 7&2^5\\ 8&2^6\\ 9&2^7\\ 10&2^8\\ 11&2^9-1\\ 12&2^{10}-2\\ 13&2^{11}-1\cdot 5\\ 14&2^{12}-2\cdot 6\\ 15&2^{13}-4\cdot 7\\ 16&2^{14}-8\cdot 8\\ 17&2^{15}-16\cdot 9\\ 18&2^{16}-32\cdot 10\\ 19&2^{17}-64\cdot 11\\ 20&2^{18}-128\cdot 12. \end{array}\]Therefore, there are $2^{18}- 2^9\cdot 3$ positive integers $x$ with $d(x)=20$. This factors as $2^9(2^{9}-3)=2^9(509)$, where $509$ is prime. Thus the answer is $\boxed{511}$.

 


\end{enumerate}